% ============================================================
% f2012_fujian_arts_q06_flowchart.tex
% 2012 福建卷（文科）第 6 题 程序框图（TikZ 矢量重绘）
% 开始 → k=1,s=1 → 判断 k<4
%   是：s=2s-k → k=k+1 → 回到判断
%   否：输出 s → 结束
% 来源 PDF：普通高考/2012/2012福建文.pdf，印面 731 第 1 栏。
% ============================================================
\begin{tikzpicture}[
  x=1cm, y=1cm, line width=0.6pt, font=\footnotesize, >=Stealth,
  every node/.style={inner sep=3pt},
  proc/.style={rectangle, draw, minimum width=2.7cm, minimum height=0.6cm},
  io/.style={trapezium, draw, trapezium left angle=70, trapezium right angle=110,
             minimum width=2.5cm, minimum height=0.62cm, trapezium stretches=true},
  dec/.style={diamond, draw, aspect=2.3, minimum width=2.2cm, minimum height=0.95cm},
  term/.style={rectangle, rounded corners=8pt, draw, minimum width=1.7cm, minimum height=0.56cm}
]

  \node[term] (start) at (0,0) {开始};
  \node[proc] (init)  at (0,-1.25) {$k=1$，$s=1$};
  \node[dec]  (judge) at (0,-2.95) {$k<4$?};

  \node[proc] (sbox)  at (3.05,-2.35) {$s=2s-k$};
  \node[proc] (kbox)  at (3.05,-1.45) {$k=k+1$};

  \node[io]   (out)   at (0,-4.55) {输出 $s$};
  \node[term] (end)   at (0,-5.65) {结束};

  \draw[->] (start) -- (init);
  \draw[->] (init) -- (judge);
  \draw[->] (judge.east) -- (3.05,-2.95) -- (sbox.south);
  \draw[->] (sbox.north) -- (kbox.south);
  \draw[->] (kbox.north) -- (3.05,-0.92) -- (0,-0.92);
  \draw[->] (judge.south) -- (out.north);
  \draw[->] (out) -- (end);

  \node at (1.62,-2.74) {是};
  \node at (0.38,-3.72) {否};

\end{tikzpicture}
